\documentclass[aps,prb,twocolumn,superscriptaddress]{revtex4-2}

\usepackage{amsmath,amssymb,amsfonts}
\usepackage{graphicx}
\usepackage{hyperref}

\begin{document}

\title{Supplemental Material: Second Scaling Exponent for Off-Diagonal Coherence in Power-Law Hopping Nonequilibrium Steady States}

\author{}
\date{}

\maketitle

This Supplemental Material provides (S1)~the equation of motion for the single-particle correlation matrix and the numerical solution method; (S2)~the proof that off-diagonal correlations are purely imaginary for power-law hopping; (S3)~the leading scaling ansatz for the $1/\alpha$ term in $\beta$; (S4)~the model-selection analysis supporting the dephasing-plateau leading law; (S5)~the first-order commutator expansion and the correction factor $\kappa$; and (S6)~tabulated data, finite-size stability checks, and code-to-data correspondence. Equation, figure, and reference numbers prefixed with ``S'' refer to this document; unprefixed numbers refer to the main text.

%=====================================================================
\section{Equation of motion and numerical method}
\label{sec:S1}
%=====================================================================

The model is a one-dimensional chain of $L$ spinless free fermions with power-law hopping $h_{ij} = J_0|i-j|^{-\alpha}$ ($J_0 = 0.3$), boundary driving at sites $1$ and $L$, and bulk dephasing at rate $\gamma_\phi$ on every site. Because the Hamiltonian is quadratic and all jump operators are linear (boundary) or bilinear-diagonal (dephasing) in the fermion operators, the steady state is Gaussian and is completely specified by the single-particle correlation matrix $C_{ij} = \operatorname{Tr}[c_i^\dagger c_j\,\rho_{\rm ss}]$.

The Lindblad equation for $\rho$ implies a closed linear equation of motion for $C$:
\begin{equation}
\frac{dC}{dt} = i(h^T C - C h^T) - \tfrac{1}{2}\{\Gamma, C\} - \gamma_\phi (C - \operatorname{diag}(C)) + \Gamma^{\text{in}} = 0 \quad (\text{steady state}),
\tag{S1}
\end{equation}
where $\Gamma$ is the diagonal matrix of total boundary rates ($\Gamma_{11} = \Gamma_L$, $\Gamma_{LL} = \Gamma_R$, zero otherwise), $\Gamma^{\text{in}} = \operatorname{diag}(\Gamma_L f_L, 0, \dots, 0, \Gamma_R f_R)$ is the injection source, and the dephasing term $-\gamma_\phi(C - \operatorname{diag} C)$ damps every off-diagonal element at rate $\gamma_\phi$ while leaving the diagonal (occupations) untouched.

\textbf{Solution method.} Equation (S1) is a Sylvester-type linear system in the $L^2$ unknowns $C_{ij}$. We vectorize $C \to \operatorname{vec}(C)$ and assemble the $L^2 \times L^2$ coefficient matrix $A$ such that $A\cdot\operatorname{vec}(C) = \operatorname{vec}(\Gamma^{\text{in}})$. For the baseline scan $L = 4, 8, 16, 32, 64$, the dimension $L^2 = 16\ldots4096$ is well within dense-solver range and we use dense LU factorization (\texttt{numpy.linalg.solve}). For the $L=128$ and $L=256$ finite-size scaling anchors, we use the sparse site-basis GMRES implementation in \texttt{fss\_compute.py}, with diagonal preconditioning and residual tolerance $\text{rtol} = 10^{-10}$. The dense baseline solver is implemented in \texttt{phase2\_L64\_compute.py} / \texttt{phase2\_L\_large.py} (COH scan) and \texttt{firewall\_AHA1.py} (universality scan).

\textbf{Cross-validation.} As an independent check, \texttt{defense\_all\_attacks.py} re-solves (S1) with three additional routes: (i)~\texttt{scipy.linalg.solve\_continuous\_lyapunov}, (ii)~an eigenbasis Green's-function evaluation [Eq.~(S3) below], and (iii)~Lyapunov iteration to $L = 128$. All three agree with the dense site-basis solver to relative error ${<}10^{-8}$ on $|C_{\rm mid}|$, and the condition number of $A$ stays below $10^5$ for all reported $(L, \alpha, \gamma_\phi)$.

%=====================================================================
\section{Pure-imaginary off-diagonals: proof for power-law hopping}
\label{sec:S2}
%=====================================================================

We prove that $\operatorname{Re}(C_{ij}) = 0$ for all $i \neq j$, generalizing the nearest-neighbor result of Bhat and \v{Z}nidari\v{c} [main-text Ref.~8] to arbitrary power-law hopping. The proof does not rely on transfer matrix methods (which require tridiagonal structure) and instead uses the structure of the steady-state Lyapunov equation directly.

The single-particle correlation matrix $C$ satisfies the steady-state equation (S1). Separate $C = D + O$ where $D = \operatorname{diag}(C)$ (real occupations) and $O$ contains the off-diagonal elements. The real part of the off-diagonal block of (S1) reads:
\begin{equation}
i[h, \operatorname{Re}(O)] - \gamma_\phi \operatorname{Re}(O) - \tfrac{1}{2}\{\Gamma, \operatorname{Re}(O)\} = 0 \quad (i \neq j).
\tag{S2}
\end{equation}

Crucially, the injection source $\Gamma^{\text{in}} = \operatorname{diag}(\Gamma_L f_L, 0, \dots, 0, \Gamma_R f_R)$ is purely diagonal and real, so it does not enter the off-diagonal equations. Equation (S2) is a homogeneous linear system for the $L(L-1)$ unknowns $\operatorname{Re}(O_{ij})$, $i \neq j$. The matrix of this linear system contains the dephasing term $-\gamma_\phi$ on every diagonal element. Since $h$ is real symmetric, the commutator $i[h, \cdot]$ maps real symmetric off-diagonal matrices to real symmetric off-diagonal matrices. The total coefficient matrix acting on $\operatorname{vec}(\operatorname{Re}(O))$ for $i \neq j$ is therefore of the form $-\gamma_\phi I + M$, where $M$ encodes boundary damping and the commutator. For $\gamma_\phi > 0$ the matrix $-\gamma_\phi I + M$ is nonsingular, and the unique solution of the homogeneous system is $\operatorname{Re}(O_{i\neq j}) = 0$. This argument uses only (a)~$h$ is real symmetric, (b)~the source is diagonal and real, and (c)~$\gamma_\phi > 0$ --- all three insensitive to the hopping range. $\square$

\textbf{Numerical confirmation.} Across all 125 COH points and 60 universality points, $\max|\operatorname{Re}(C_{i\neq j})| < 10^{-15}$ (e.g., the firewall purity check reports $\max_{\operatorname{Re}} = 3.7\times10^{-17}$ against $\max_{\operatorname{Im}} = 4.5\times10^{-2}$, ratio $8\times10^{-16}$). The property is structural, not the spatial-decay phenomenon studied in the main text.

%=====================================================================
\section{Leading scaling ansatz for $\beta(\alpha) \propto 1/\alpha$}
\label{sec:S3}
%=====================================================================

This section gives the constrained leading-scaling ansatz for the coherence-decay exponent. The exact object is not a closed-form theorem for $\beta$ at all $\gamma_\phi$ and all finite $L$; it is a leading scaling structure selected by the fractional kernel, two fixed-point limits, and a collapse/model-selection test. Steps S3.1--S3.5 establish the reduction to a fractional-boundary scaling problem. Steps S3.6--S3.7 define the leading ansatz and its correction windows.

\subsection*{Step 1 (rigorous) --- Off-diagonal slaved to occupation gradient}

A first-order expansion of (S1) in $\gamma_\phi^{-1}$ [derived in S5] gives, for the midchain nearest-neighbor pair,
\begin{equation}
|C_{\rm mid}| = \frac{J_0}{\gamma_\phi}\,|D_{L/2+1} - D_{L/2}|\,\kappa(L,\alpha),
\tag{S3}
\end{equation}
with $\kappa \to 1$ as $L \to \infty$. Hence in the thermodynamic limit $\beta(\alpha,\gamma_\phi) = \nu(\alpha,\gamma_\phi)$, where $|D_{L/2+1} - D_{L/2}| \sim L^{-\nu}$.

\subsection*{Step 2 (rigorous) --- Interior occupations solve a discrete fractional Laplacian}

Setting the time derivative of $C_{ii}$ to zero in the bulk gives $\sum_{r\neq 0}(D_{i+r} - D_i)/r^{2\alpha} = 0$. The kernel $r^{-2\alpha}$ defines a nonlocal operator $\mathcal{L}_\alpha$.

\subsection*{Step 3 (rigorous) --- Fourier kernel and its small-$k$ behavior}

With $\mathcal{L}_\alpha(k) = -2\sum_{r\geq 1}(1 - \cos kr)/r^{2\alpha}$,
\begin{equation}
\mathcal{L}_\alpha(k) \approx
\begin{cases}
-D_{\rm eff}\,k^2, & \alpha > 3/2 \quad\text{(normal diffusion, $D_{\rm eff} = \sum_r r^{2-2\alpha} < \infty$)}; \\[4pt]
-c_\alpha|k|^{2\alpha-1}, & \alpha < 3/2 \quad\text{(anomalous diffusion)}.
\end{cases}
\tag{S4}
\end{equation}
The crossover at $\alpha = 3/2$ is where $\sum_r r^{2-2\alpha} = \sum_r r^{-1}$ diverges logarithmically --- the analytic origin of the ballistic-to-diffusive transition of Dhawan et al.~[Ref.~5].

\subsection*{Step 4 (rigorous) --- Nonlocal Robin boundary conditions}

The boundary occupation obeys
\begin{equation}
\frac{2J_0^2}{\gamma_\phi}\sum_{k>1}\frac{D_k - D_1}{k^{2\alpha}} + \Gamma_L(D_1 - f_L) = 0.
\tag{S5}
\end{equation}
The nonlocal sum carries a prefactor $\sim L^{1-2\alpha}$, which diverges for $\alpha < 3/2$: the boundary condition couples to all interior sites, not just the nearest one.

\subsection*{Step 5 (rigorous) --- Boundary-layer width}

Rescaling $x = y/L$, the competition between the $L^{1-2\alpha}$ nonlocal prefactor and the $|k|^{2\alpha-1}$ bulk dispersion produces a boundary layer of width $\delta(\alpha,L) \sim L^{(2\alpha-1)/(3-2\alpha)}$ for $\alpha < 3/2$. As $\alpha \to 1^+$, $\delta \sim L$ (fully nonlocal); as $\alpha \to 3/2^-$, $\delta$ diverges (crossover signal).

\subsection*{Step 6 (fixed-point constraints) --- Midpoint gradient}

The midchain gradient inherits the boundary-layer profile, $|D'(L/2)| \sim (\Delta f/L)\cdot F'_\alpha(1/2; L)$, where $F_\alpha$ is the fractional harmonic profile. The effective exponent is $\nu(\alpha) = 1 - d\ln F'_\alpha/d\ln L$. Two fixed points constrain the leading form. In the ballistic/dephasing-irrelevant limit $\gamma_\phi \to 0$, the steady-state gradient vanishes and the $\alpha$-dependent amplitude must collapse, $a(\gamma_\phi) \to 0$. In the strong-dephasing or short-range limit, the nonlocal Robin layer flows toward the ordinary diffusive boundary problem, so $\beta \to 1$ and the long-range correction is suppressed. Thus the exponent can be written as $\beta = b(\gamma_\phi) + a(\gamma_\phi)X(\alpha) + \delta\beta_{\rm corr}$, where $X$ is the long-range scaling coordinate and $\delta\beta_{\rm corr}$ collects boundary-layer and finite-$L$ corrections.

\subsection*{Step 7 (leading scaling coordinate + calibration) --- $\beta$ functional form}

For power-law hopping $h(r) \sim r^{-\alpha}$, the nonlocal dispersion and Robin boundary prefactor depend on $\alpha$ through inverse powers of the hopping decay exponent. On compact $\alpha$ intervals away from the singular boundary-layer edge $\alpha = 1$, the leading long-range coordinate is $X(\alpha) = 1/\alpha$, with higher analytic and finite-size corrections absorbed into $\delta\beta_{\rm corr} = c_1/\alpha^2 + c_2 L^{-\omega} + \cdots$. Combining Steps~1 and~6 gives
\begin{equation}
\beta(\alpha,\gamma_\phi) = \frac{a(\gamma_\phi)}{\alpha} + b(\gamma_\phi) + \delta\beta_{\rm corr}(\alpha,\gamma_\phi,L).
\tag{S6}
\end{equation}

The $\gamma_\phi$-dependence of the coefficients follows from the hopping-dephasing competition: $a(\gamma_\phi)$ peaks where $J_0$ and $\gamma_\phi$ compete and vanishes toward the ballistic fixed point; $b(\gamma_\phi)$ rises toward the diffusive baseline. The numerical values of $a$ and $b$ are calibrated in S4. The falsifiable prediction is a collapse test: in the leading scaling window, $[\beta(\alpha,\gamma_\phi) - b(\gamma_\phi)]/a(\gamma_\phi)$ should be linear in $1/\alpha$, while the correction-dominated windows $\gamma_\phi = 0.01$, $\gamma_\phi = 0.1$, and $\alpha \approx 1.1$ should show the largest controlled deviations. \textit{[Leading scaling ansatz; amplitudes and corrections calibrated.]}

\textbf{Limit checks.}
\begin{itemize}
\item $\gamma_\phi \to 0$: no steady-state gradient $\Rightarrow$ $a(\gamma_\phi) \to 0$ and $\beta$ becomes nearly $\alpha$-independent. Data: $\beta \approx 0.15$ at $\gamma_\phi = 0.01$, consistent with a ballistic/correction-dominated window.
\item $\gamma_\phi \to \infty$: pure diffusion $\Rightarrow$ $\beta \to 1$ after boundary corrections vanish. Data: $\beta \approx 1.0$--$1.23$ at $\gamma_\phi = 2.0$, consistent with an approach to the diffusive baseline plus finite-$L$ boundary corrections.
\item $\alpha \to 1^+$: $\beta = a + b = 1.235$ ($\gamma_\phi = 0.5$); data $\beta(1.1) = 1.172$, extrapolation error $\sim 5\%$.
\item $\alpha \to \infty$: $\beta = b = 0.422$ is a finite-$L$ ($L \leq 64$) extrapolated value; the asymptotic $\beta \to 1$ requires $L \gtrsim 256$ and is flagged as unverified in the main text.
\end{itemize}

%=====================================================================
\section{Model selection for the $1/\alpha$ functional form}
\label{sec:S4}
%=====================================================================

For each $\gamma_\phi$ we fit the five $\beta(\alpha)$ values ($\alpha = 1.1, \dots, 1.9$) to four candidate forms by least squares and rank them by the small-sample Akaike information criterion $\text{AIC}_c = n\ln(\text{RSS}/n) + 2k + 2k(k+1)/(n-k-1)$, with $n = 5$ and $k$ the number of fitted parameters. Computed by \texttt{model\_selection.py}. Model selection is used here to identify the leading dephasing-plateau law and the correction-dominated windows; it is not used to claim that one two-parameter curve is globally optimal at every $\gamma_\phi$.

\begin{table}[h]
\caption{Model selection at $\gamma_\phi = 0.5$ (lower $\text{AIC}_c$ is better).}
\label{tab:S1}
\begin{ruledtabular}
\begin{tabular}{lcccccc}
Model & params ($k$) & fitted coefficients & RSS & $R^2$ & RMSE & $\text{AIC}_c$ \\
\hline
$\beta = a/\alpha + b$ & 2 & $a=0.8127$, $b=0.4218$ & $1.15\times10^{-3}$ & 0.981 & 0.0152 & $\mathbf{-31.9}$ \\
$\beta = a\ln\alpha + b$ & 2 & $a=-0.563$, $b=1.202$ & $2.77\times10^{-3}$ & 0.955 & 0.0235 & $-27.5$ \\
$\beta = m\alpha + c$ & 2 & $m=-0.377$, $c=1.549$ & $5.16\times10^{-3}$ & 0.917 & 0.0321 & $-24.4$ \\
$\beta = a/\alpha^2 + b/\alpha + c$ & 3 & $a=0.881$, $b=-0.454$, $c=0.861$ & $2.47\times10^{-4}$ & 0.996 & 0.0070 & $-19.6$ \\
\end{tabular}
\end{ruledtabular}
\end{table}

At $\gamma_\phi = 0.5$, the two-parameter $1/\alpha$ form has the lowest $\text{AIC}_c$ for the $L\leq 64$ baseline table. The three-parameter form reduces RSS by $\sim 5\times$ but is penalized by the small-sample term: at $n = 5$, adding a parameter costs heavily, and $\text{AIC}_c$ rises to $-19.6$. Thus $\gamma_\phi = 0.5$ is the clean baseline point where the leading $1/\alpha$ law is selected over the two-parameter alternatives. The $L=128$ and $L=256$ anchors in Table~S7 show that this baseline law is the small-to-intermediate $\alpha$ branch of a larger finite-size flow, not a global asymptotic curve.

\begin{table}[h]
\caption{Calibrated $(a, b)$ and fit quality across $\gamma_\phi$ (from \texttt{phase2\_fit\_results\_FULL.json} regressions).}
\label{tab:S2}
\begin{ruledtabular}
\begin{tabular}{cccccc}
$\gamma_\phi$ & $a(\gamma_\phi)$ & $b(\gamma_\phi)$ & $R^2$ ($1/\alpha$ fit) & note \\
\hline
0.01 & 0.058 & 0.115 & 0.996${}^*$ & ballistic; $1/\alpha$ form physically inapplicable \\
0.1  & 0.357 & 0.425 & 0.960 & weak dephasing; linear model $\text{AIC}_c = -50.53$ (lower than $1/\alpha$) \\
0.5  & 0.813 & 0.422 & 0.981 & optimal competition regime; $1/\alpha$ selected by $\text{AIC}_c$ \\
1.0  & 0.745 & 0.529 & 0.949 & strong dephasing \\
2.0  & 0.610 & 0.657 & 0.940 & near-diffusive \\
\end{tabular}
\end{ruledtabular}
${}^*$At $\gamma_\phi = 0.01$ the linear-in-$1/\alpha$ regression is statistically good but physically vacuous: the five $\beta$ values are themselves $1\sigma$-consistent with a constant (${\sim}0.15$), so the apparent $a$, $b$ carry no scaling information.
\end{table}

\textbf{Full model selection across $\gamma_\phi$.} Table~\ref{tab:S2b} reports $\text{AIC}_c$ for all four candidate forms at each $\gamma_\phi$ ($n=5$ throughout; computed in \texttt{model\_selection.py}). Lower $\text{AIC}_c$ is better. At $\gamma_\phi = 0.5$, $1/\alpha$ is the best two-parameter model; at $\gamma_\phi = 0.1$, the linear model wins ($\Delta\text{AIC}_c > 14$). The $1/\alpha$ form is therefore the leading long-range coordinate in the dephasing window where hopping and dephasing compete, not a global functional form at every $\gamma_\phi$.

\begin{table}[h]
\caption{Full model selection: $\text{AIC}_c$ for all candidate forms at each $\gamma_\phi$.}
\label{tab:S2b}
\begin{ruledtabular}
\begin{tabular}{cccccc}
$\gamma_\phi$ & $1/\alpha$ ($k=2$) & $\ln\alpha$ ($k=2$) & linear ($k=2$) & $1/\alpha^2 + 1/\alpha + c$ ($k=3$) \\
\hline
0.01 & $-66.12$ & $-78.34$ & $-61.51$ & $-77.58$ \\
0.1  & $-36.14$ & $-41.13$ & $\mathbf{-50.53}$ & $-43.04$ \\
0.5  & $\mathbf{-31.89}$ & $-27.50$ & $-24.38$ & $-19.57$ \\
1.0  & $\mathbf{-27.52}$ & $-24.58$ & $-22.36$ & $-29.17$ \\
2.0  & $-28.73$ & $-25.98$ & $-23.88$ & $\mathbf{-34.13}$ \\
\end{tabular}
\end{ruledtabular}
\end{table}

At $n=5$ the small-sample $\text{AIC}_c$ penalty for $k=3$ is severe ($2k(k+1)/(n-k-1) = 24$), so the two-parameter vs.\ three-parameter comparison should be interpreted with caution. The two-parameter comparisons ($1/\alpha$ vs.\ $\ln\alpha$ vs.\ linear) are the statistically meaningful ones at this sample size.

%=====================================================================
\section{First-order commutator expansion and the correction factor $\kappa$}
\label{sec:S5}
%=====================================================================

Write $C = D + O$ with $D = \operatorname{diag}(C)$ (occupations) and $O$ the off-diagonal part. The off-diagonal block of the steady-state equation (S1) reads $i[h, D + O] - \gamma_\phi O = 0$ (boundary terms are subleading in the bulk). To leading order in $\gamma_\phi^{-1}$, the commutator $[h, O]$ is higher order and
\begin{equation}
O_{ij} = \frac{i\,(h D - D h)_{ij}}{\gamma_\phi} = \frac{i\,h_{ij}(D_j - D_i)}{\gamma_\phi} + \mathcal{O}(\gamma_\phi^{-3}).
\tag{S7}
\end{equation}

This shows directly that $O$ is purely imaginary (consistent with S2) and that, for the midchain nearest-neighbor pair, $|C_{\rm mid}^{(1)}| = (J_0/\gamma_\phi)|D_{L/2+1} - D_{L/2}|$. The correction factor $\kappa \equiv |C_{\rm exact}|/|C_{\rm mid}^{(1)}|$ measures the higher-order commutator contribution.

\begin{table}[h]
\caption{$\kappa$ at $\gamma_\phi = 0.5$, computed from the exact $|C_{\rm mid}|$ and the stored occupation profile \texttt{diag\_C} (script \texttt{model\_selection.py}, output \texttt{model\_selection\_results.json}).}
\label{tab:S3}
\begin{ruledtabular}
\begin{tabular}{ccccc}
$\alpha$ \textbackslash $L$ & 16 & 32 & 64 \\
\hline
1.1 & 1.073 & 1.035 & 1.017 \\
1.3 & 1.049 & 1.019 & 1.006 \\
1.5 & 1.030 & 1.008 & 1.002 \\
1.7 & 1.017 & 1.003 & 1.000 \\
1.9 & 1.009 & 1.001 & 1.000 \\
\end{tabular}
\end{ruledtabular}
\end{table}

$\kappa \to 1$ monotonically with $L$ for every $\alpha$; the slowest convergence is the longest-range case $\alpha = 1.1$. For $L \geq 32$, $\kappa < 1.04$; for $L \geq 64$, $\kappa < 1.02$. This justifies the identification $\beta = \nu$ in the thermodynamic limit (Step~1 of S3).

%=====================================================================
\section{Data tables, convergence, and code-to-data correspondence}
\label{sec:S6}
%=====================================================================

\begin{table}[h]
\caption{Full $\beta(\alpha, \gamma_\phi)$ with $1\sigma$ standard errors from log-log regression of $|C_{\rm mid}|$ vs.\ $L$ over $L = 4, 8, 16, 32, 64$ (reproduces main-text Table~1; source \texttt{phase2\_fit\_results\_FULL.json}).}
\label{tab:S4}
\begin{ruledtabular}
\begin{tabular}{ccccccc}
$\alpha$ \textbackslash $\gamma_\phi$ & 0.01 & 0.1 & 0.5 & 1.0 & 2.0 \\
\hline
1.1 & 0.167(37) & 0.739(65) & 1.172(11) & 1.230(32) & 1.234(35) \\
1.3 & 0.160(34) & 0.710(63) & 1.043(14) & 1.080(32) & 1.106(34) \\
1.5 & 0.154(32) & 0.675(58) & 0.943(15) & 0.997(26) & 1.040(29) \\
1.7 & 0.149(30) & 0.636(51) & 0.891(8)  & 0.962(18) & 1.012(22) \\
1.9 & 0.145(29) & 0.602(46) & 0.872(3)  & 0.951(10) & 1.004(17) \\
\end{tabular}
\end{ruledtabular}
(The $\gamma_\phi = 0.1$ and $\gamma_\phi \geq 1.0$ standard errors reported here are the raw regression std\_err from the data files; main-text Table~1 quotes a subset. $R^2 > 0.997$ for $\gamma_\phi \geq 0.5$; $R^2 \approx 0.96$--$0.98$ for $\gamma_\phi = 0.1$; $R^2 \approx 0.87$--$0.89$ for $\gamma_\phi = 0.01$.)
\end{table}

\begin{table}[h]
\caption{Anti-correlation slopes $d\beta/d\mu$ (full-range linear fit over $\alpha \in [1.1, 1.9]$, $\mu = 2\alpha - 2$; source \texttt{model\_selection\_results.json}).}
\label{tab:S5}
\begin{ruledtabular}
\begin{tabular}{cccc}
$\gamma_\phi$ & $d\beta/d\mu$ & $R^2$ \\
\hline
0.01 & $-0.014$ & 0.990 \\
0.1  & $-0.087$ & 0.998 \\
0.5  & $-0.188$ & 0.917 \\
1.0  & $-0.170$ & 0.856 \\
2.0  & $-0.138$ & 0.843 \\
\end{tabular}
\end{ruledtabular}
\end{table}

The slope is strictly negative for all $\gamma_\phi \geq 0.1$ and varies by a factor ${>}2$, confirming the anti-correlation is $\gamma_\phi$-dependent and hence not a fixed algebraic relation in $\alpha$ alone.

\begin{table}[h]
\caption{Finite-size stability of $\beta$ at the reference dephasing point $\gamma_\phi = 0.5$. $\beta_{\rm all}$ uses $L = 4, 8, 16, 32, 64$; $\beta_{\rm drop\,L4}$ uses $L = 8, 16, 32, 64$; $\beta_{L\geq 16}$ uses $L = 16, 32, 64$. The leave-one-size-out column reports the maximum absolute shift in $\beta$ after deleting any one system size. Source: \texttt{finite\_size\_stability.py} / \texttt{finite\_size\_stability\_results.json}.}
\label{tab:S6}
\begin{ruledtabular}
\begin{tabular}{cccccc}
$\alpha$ & $\beta_{\rm all}$ & $\beta_{\rm drop\,L4}$ & $\beta_{L\geq 16}$ & max leave-one-size shift & drop-L4 shift \\
\hline
1.1 & 1.1723 & 1.1544 & 1.1482 & 0.0179 & 1.52\% \\
1.3 & 1.0428 & 1.0220 & 1.0006 & 0.0216 & 1.99\% \\
1.5 & 0.9427 & 0.9213 & 0.8982 & 0.0213 & 2.26\% \\
1.7 & 0.8912 & 0.8786 & 0.8685 & 0.0127 & 1.42\% \\
1.9 & 0.8715 & 0.8708 & 0.8758 & 0.0043 & 0.09\% \\
\end{tabular}
\end{ruledtabular}
\end{table}

The leading $\alpha$-dependence in the dephasing plateau is therefore not driven by the smallest system. Deleting $L = 4$ changes $\beta$ by at most 2.26\% at $\gamma_\phi = 0.5$. At $\gamma_\phi = 0.1$, the $L=4$ deletion shifts are larger (12--13\%, from \texttt{finite\_size\_stability\_results.json}) --- consistent with the weaker dephasing making the smallest system a poorer approximation to the thermodynamic scaling regime. These shifts are the reason we treat $\gamma_\phi = 0.5$ as the reference dephasing window and $\gamma_\phi = 0.1$ as a crossover regime. The larger $L=128$ and $L=256$ anchors below further reveal that the asymptotic flow bends into a near-diffusive plateau rather than remaining globally monotone.

\begin{table}[h]
\caption{$L=256$ finite-size scaling anchors at $\gamma_\phi = 0.5$. $\beta(L{\geq}32)$ and $\beta(L{\geq}64)$ are log-log fits using $L$ up to 256. $\beta_{64\to128}$ and $\beta_{128\to256}$ are local two-size exponents. Source: \texttt{fss\_gamma0.5\_L256.json} / \texttt{fss\_gamma0.5\_L256\_raw.json}.}
\label{tab:S7}
\begin{ruledtabular}
\begin{tabular}{ccccccc}
$\alpha$ & $|C_{\rm mid}|(L{=}256)$ & $\beta(L{\geq}32)$ & $\beta(L{\geq}64)$ & $\beta_{64\to128}$ & $\beta_{128\to256}$ \\
\hline
1.1 & $1.880190\times10^{-4}$ & 1.1162 & 1.1045 & 1.1160 & 1.0929 \\
1.3 & $3.748596\times10^{-4}$ & 0.9472 & 0.9339 & 0.9415 & 0.9262 \\
1.5 & $5.549012\times10^{-4}$ & 0.8901 & 0.8952 & 0.8850 & 0.9054 \\
1.7 & $6.445630\times10^{-4}$ & 0.9040 & 0.9198 & 0.9045 & 0.9351 \\
1.9 & $6.777337\times10^{-4}$ & 0.9291 & 0.9463 & 0.9326 & 0.9599 \\
\end{tabular}
\end{ruledtabular}
\end{table}

The $L=256$ anchors show a minimum near $\alpha \approx 1.5$ and a large-$\alpha$ return toward a near-diffusive plateau. This structure is reproducible in the local exponents and is therefore treated as part of the physical finite-size flow, not as a fitting artifact.

\textbf{Finite-size note.} The $L{\leq}64$ $\beta$ table is the baseline scan. The $L=128$ slices and the $L=256$ $\gamma_\phi = 0.5$ anchors show that the large-$L$ coherence exponent should be interpreted as a flow: the small-$\alpha$ branch remains above $\beta \approx 1$, while the large-$\alpha$ branch approaches a near-diffusive plateau. The remaining open finite-size problem is not whether midchain $\beta$ exists at $L=128$, but how far the plateau values drift beyond $L=256$ and how the $\gamma_\phi = 0.01$ ballistic limit connects to this flow.

\subsection*{Code-to-data correspondence}

\begin{table}[h]
\caption{Code-to-data correspondence.}
\label{tab:S8}
\begin{ruledtabular}
\begin{tabular}{lll}
Script & Produces & Used for \\
\hline
\texttt{phase2\_L64\_compute.py}, \texttt{phase2\_L\_large.py} & \texttt{phase2\_raw\_results\_FULL.json}, \texttt{phase2\_fit\_results\_FULL.json} & $\beta(\alpha,\gamma_\phi)$, Tables~1/S2/S4, Figs~1--2 \\
\texttt{firewall\_AHA1.py} & \texttt{firewall\_AHA1\_results.json} & universality test, Table~3, Fig~3 \\
\texttt{model\_selection.py} & \texttt{model\_selection\_results.json} & $\text{AIC}_c$ (Table~S1), $d\beta/d\mu$ (Table~S5), $\kappa$ (Table~S3) \\
\texttt{finite\_size\_stability.py} & \texttt{finite\_size\_stability\_results.json} & finite-size stability checks, Table~S6 \\
\texttt{fss\_compute.py} & \texttt{fss\_gamma*.json}, \texttt{fss\_gamma0.5\_L256.json} & $L{=}128$/$L{=}256$ FSS flow, Table~S7 \\
\texttt{defense\_all\_attacks.py}, \texttt{firewall1\_L128\_v2.py} & $L{=}128$ + cross-validation & convergence / solver checks \\
\texttt{figures/generate\_figures.py} & Fig1--3.\{pdf,png\} & main-text figures \\
\end{tabular}
\end{ruledtabular}
\end{table}

All numerical results in the main text and this Supplement are reproducible by running these scripts on the archived data (Zenodo DOI: \texttt{10.5281/zenodo.20553343}); no value is hand-entered or model-generated.

%=====================================================================
% Bibliography (shared with main text)
%=====================================================================
\begin{thebibliography}{8}

\bibitem{breuer}
H.-P.~Breuer and F.~Petruccione,
\textit{The Theory of Open Quantum Systems}
(Oxford University Press, Oxford, 2002).

\bibitem{landi}
G.~T.~Landi, D.~Poletti, and G.~Schaller,
Rev.\ Mod.\ Phys.\ \textbf{94}, 045006 (2022).

\bibitem{dhar}
A.~Dhar,
Adv.\ Phys.\ \textbf{57}, 457 (2008).

\bibitem{lepri}
S.~Lepri, R.~Livi, and A.~Politi,
Phys.\ Rep.\ \textbf{377}, 1 (2003).

\bibitem{dhawan}
A.~Dhawan, K.~Ganguly, M.~Kulkarni, and B.~K.~Agarwalla,
Phys.\ Rev.\ B \textbf{110}, L081403 (2024).

\bibitem{costa}
J.~Costa, P.~Ribeiro, and A.~De Luca,
arXiv:2504.00188v3 (2025).

\bibitem{sarkar}
S.~Sarkar, B.~K.~Agarwalla, and D.~S.~Bhakuni,
Phys.\ Rev.\ B \textbf{109}, 165408 (2024).

\bibitem{bhat}
J.~M.~Bhat and M.~\v{Z}nidari\v{c},
Phys.\ Rev.\ B \textbf{111}, 174306 (2025).

\bibitem{compute_mu}
See \texttt{compute\_mu.py} and \texttt{compute\_mu\_results.json}
in the project repository.

\end{thebibliography}

% ======================================================================
\section{Self-Consistent Transport Exponent $\mu(\alpha,\gamma_\phi)$}
\label{sec:S7}
% ======================================================================

The transport exponent $\mu$ is defined via the steady-state current scaling $J(L) \sim L^{-\mu}$. In the boundary-driven Lindblad setup, the current is spatially uniform and can be evaluated at the boundary:
\begin{equation}
J = \Gamma_L(f_L - D_1),
\label{eq:S8}
\end{equation}
where $D_1 = \langle c_1^\dagger c_1 \rangle = C_{11}$ (real) is the occupation of the first site. This follows directly from the diagonal Lindblad equation for the occupation number,
\begin{equation}
\frac{d}{dt}\langle n_1 \rangle = \Gamma_L f_L(1 - \langle n_1\rangle) - \Gamma_L(1-f_L)\langle n_1\rangle + \text{(bulk commutator)} = 0 \quad\text{(steady state)},
\label{eq:S9}
\end{equation}
giving $J = \Gamma_L(f_L - D_1)$. The same current can be obtained from the right boundary: $J = \Gamma_R(D_L - f_R)$, providing an internal consistency check. In our computations, the two boundary estimates agree to within machine precision.

We compute $J(L)$ for $L = 8, 16, 32, 64$ (and $L = 128$ for $\gamma_\phi = 0.5$) and extract $\mu$ from log-log linear regression. Full results are stored in \texttt{compute\_mu\_results.json}.

\begin{table}[h]
\caption{Self-consistent transport exponent $\mu(\alpha, \gamma_\phi)$ with $1\sigma$ standard errors from log-log regression of $J$ vs.\ $L$ ($L = 8, 16, 32, 64$; $L$ up to 128 for $\gamma_\phi=0.5$). All fits have $R^2 > 0.99$.}
\label{tab:S8}
\begin{ruledtabular}
\begin{tabular}{c|cccc}
$\alpha$ & $\mu(\gamma_\phi{=}0.1)$ & $\mu(\gamma_\phi{=}0.5)$ & $\mu(\gamma_\phi{=}1.0)$ & $\mu(\gamma_\phi{=}2.0)$ \\
\hline
$1.1$ & $0.2696(37)$ & $0.2846(92)$ & $0.3032(89)$ & $0.3122(90)$ \\
$1.3$ & $0.3144(87)$ & $0.4146(22)$ & $0.4608(9)$  & $0.5026(14)$ \\
$1.5$ & $0.3664(161)$ & $0.5792(173)$ & $0.6301(117)$ & $0.6833(76)$ \\
$1.7$ & $0.4237(261)$ & $0.7178(258)$ & $0.7644(174)$ & $0.8173(110)$ \\
$1.9$ & $0.4792(361)$ & $0.8071(264)$ & $0.8519(173)$ & $0.9010(97)$ \\
\end{tabular}
\end{ruledtabular}
\end{table}

\paragraph{Comparison with Dhawan et al. [5].}
The dephasing-free formula $\mu(\alpha) = 2\alpha-2$ (valid for $1 < \alpha < 3/2$ in coherent systems) predicts $\mu(1.1) = 0.2$, $\mu(1.3) = 0.6$, $\mu(1.5) = 1.0$. Our computed $\mu$ values are systematically lower, and the discrepancy grows with $\gamma_\phi$. This is physically expected: bulk dephasing introduces an additional dissipation channel that suppresses the current relative to the coherent case. At $\alpha = 1.5$, the Dhawan formula predicts diffusive transport ($\mu = 1.0$), while our model with $\gamma_\phi = 0.5$ gives $\mu = 0.579$---the system remains subdiffusive because dephasing randomizes the phase coherence needed for efficient transport. The $\gamma_\phi$-sensitivity of $\mu$ confirms that the transport exponent in the presence of bulk dephasing is not universal in the same sense as the dephasing-free case.

\end{document}
